\section{Convergence and constructive multiscale closure}\label{bmh:sec:closure}

\subsection{The product strip}
\begin{proposition}[Full common convergence strip]\label{bmh:prop:strip}
For $r$ polylogarithm factors with scales in
$\C\setminus(-\infty,0]$, the ordinary integral
\[
 \bmhM_N(a;\boldsymbol s;\boldsymbol z)=
 \int_0^\infty\frac{x^{a-1}}{(1+x)^N}
       \prod_{j=1}^r\Li_{s_j}(-z_jx)\dd x
\]
is holomorphic for $-r<\Rea a<N$ and entire in each spectral order.
Arbitrary fixed spectral derivatives and powers of $\log x$ can be
inserted under the integral on compact subsets of this domain.
This is the maximal strip common to all spectral orders.
\end{proposition}
\begin{proof}
At zero the power series gives $\Li_s(-zx)=O(x)$ uniformly on compact
sets of $s,z$. At infinity the bound is $O((1+\log x)^L)$ for some
compact-dependent $L$. Here is a direct justification of the latter
bound that includes nonpositive spectral orders. Choose an integer $k$
such that $\Rea(s+k)>0$ on the spectral compact set. The Fermi integral
\[
 \Li_{s+k}(-zx)=-\frac1{\Gamma(s+k)}\int_0^\infty
 u^{s+k-1}\frac{zx e^{-u}}{1+zx e^{-u}}\dd u
\]
can be differentiated $k$ times by $x\partial_x$. The differentiated
rational kernels are bounded by a constant times
$\min(xe^{-u},1)$ on the chosen scale compact. Split the $u$ integral at
$\log x$. This proves a fixed power-logarithmic bound. Small circles in
the spectral variables and Cauchy's formula give the same conclusion
for fixed spectral derivatives. The resulting majorants are integrable
on every smaller strip; parameter holomorphy and differentiation follow.
Finally $s_j=0,z_j=1$ makes the product $(-x/(1+x))^r$, forcing both
strict endpoint inequalities.
\end{proof}

\begin{proposition}[Denominator recurrence]\label{bmh:prop:Nrec}
With $E_j$ lowering the $j$th spectral order by one,
\begin{equation}\label{bmh:eq:Nrec}
 \bmhM_{N+1}(a)=\frac1N\left(N-a-\sum_{j=1}^rE_j\right)\bmhM_N(a)
\end{equation}
on the common strip $-r<\Rea a<N$.
\end{proposition}
\begin{proof}
Integrate the derivative of
$x^a(1+x)^{-N}\prod_j\Li_{s_j}(-z_jx)$. Its endpoint values vanish by
the preceding estimates, and
$x\partial_x\Li_s(-zx)=\Li_{s-1}(-zx)$. Rewriting
$x/(1+x)=1-(1+x)^{-1}$ gives the formula.
\end{proof}
This recurrence is useful, but a termwise continuation beyond its common
strip can obscure cancellations. The compact-interval recurrence in
Section~\ref{bmh:sec:kernels} has no such ambiguity at integer $a$.

\subsection{A finite alphabet and a weight refinement}
Use $t=x/(1+x)$, so
\begin{equation}\label{bmh:eq:compact-kernel}
 \frac{x^{a-1}\dd x}{(1+x)^N}
 =R_{N,a}(t)\dd t,
 \qquad R_{N,a}(t)=t^{a-1}(1-t)^{N-a-1}.
\end{equation}
Let $G(\alpha_1,\ldots,\alpha_k;t)$ be the iterated integral with kernels
$\dd t/(t-\alpha)$, with the usual logarithmic value for all-zero words.
For $z\ne1$ put $\lambda_z=(1-z)^{-1}$. Direct differentiation and the
value at zero give
\begin{equation}\label{bmh:eq:Landen-G}
 \Li_m\!\left(-\frac{zt}{1-t}\right)
 =G\big((e_0-e_1)^{m-1}(e_1-e_{\lambda_z});t\big).
\end{equation}
Words and powers on the right are noncommutative, and $G$ is extended
linearly. For $z=1$ omit $e_{\lambda_z}$. Indeed the order-one identity is
$\log(1-t)-\log(1-(1-z)t)$, and higher orders follow from
$\dd\log(t/(1-t))=\dd t/t-\dd t/(t-1)$.

For fixed positive scales, let $\bmhH_{\le d}$ be the vector space over
$\Q(z_1,\ldots,z_r)$ generated by finite endpoint values of iterated
integrals in the alphabet
\begin{equation}\label{bmh:eq:alphabet}
 \{0,1,\lambda_{z_1},\ldots,\lambda_{z_r}\}
\end{equation}
of word length at most $d$, including products whose total length is
at most $d$. Endpoint logarithms are combined before taking finite
limits. This is a \emph{representation filtration}; no assertion about
independence of numerical periods is built into its definition.

\begin{theorem}[Multiscale finite closure and residue-weight law]
\label{bmh:thm:closure}
Let $s_1,\ldots,s_r$ be positive integers, $z_j>0$, $N\ge1$, and let
$a$ be an integer with $1-r\le a\le N-1$. Put $V=\sum_j s_j$.
Every logarithmic moment of $\bmhM_N$ is a finite multiple-polylogarithm
expression in the alphabet~\eqref{bmh:eq:alphabet}, and its anchored
primitive has a constructive hyperlogarithmic expression. More precisely,
for $h\ge0$ the moment belongs to $\bmhH_{\le V+h+1}$, and
\begin{equation}\label{bmh:eq:top-weight}
 \partial_a^h\bmhM_N(a)-c_{N,a}\,
 \left.\partial_b^h\bmhM_1(b)\right|_{b=0}
 \ \in\ \bmhH_{\le V+h},
\end{equation}
where
\begin{equation}\label{bmh:eq:residue}
 c_{N,a}=\begin{cases}
 0,&a\ge1,\\
 (-1)^k\binom{N+k-1}{k},&a=-k\le0.
 \end{cases}
\end{equation}
The factor parameters on both sides of~\eqref{bmh:eq:top-weight} are the
same. Coincident scales, including $z_j=1$, are interpreted by their
convergent limits.
\end{theorem}
\begin{proof}
Since $a,N$ are integers in the stated range, $R_{N,a}$ is a Laurent
polynomial, with no pole at $t=1$. Products in~\eqref{bmh:eq:Landen-G}
are expanded by shuffle, and
$\log x=\log t-\log(1-t)$ is a weight-one hyperlogarithm.
For completeness, the integration algorithm is the following finite
induction. Decompose a rational prefactor into a derivative of a rational
function and simple-pole terms. A simple pole appends one integration
letter. Integrate the rational derivative by parts; differentiating an
iterated integral removes its first letter and lowers word length by
one. Repeating terminates at weight zero. Polynomial prefactors are
handled by the same integration-by-parts step. All poles involved are
in~\eqref{bmh:eq:alphabet}, so no new letter is introduced. This is the
standard rational-hyperlogarithm mechanism~\cite[Section 2.2]{Panzer},
with the alphabet here explicitly determined.

For the sharper assertion, write
\[
 R_{N,a}(t)=\frac{c_{N,a}}t+S'(t),\qquad S(1)=0.
\]
Extracting the coefficient of $t^{-1}$ gives~\eqref{bmh:eq:residue}.
Let $P(t)$ be the polylogarithm product multiplied by
$\log^h(t/(1-t))$. Its weight is at most $V+h$. The $c_{N,a}/t$
term is precisely the second term of~\eqref{bmh:eq:top-weight} and can
raise weight by one. For the remainder,
\[
 \int_0^1S'(t)P(t)\dd t=-\int_0^1S(t)P'(t)\dd t.
\]
There are no boundary terms: at one, $S(t)=O(1-t)$ while $P$ grows at
most logarithmically; at zero, if $a=-k$, then $S=O(t^{-k})$ and
$P=O(t^r(1+|\log t|)^h)$, with $k<r$. The case $a\ge0$ is easier.
The derivative $P'$ has weight at most $V+h-1$. Rational integration
can raise it by at most one, proving~\eqref{bmh:eq:top-weight}.

These manipulations may first be performed on $[\varepsilon,1-\varepsilon]$.
The original integrability and the displayed boundary estimates determine
the combined endpoint limits, including cancellations between individual
hyperlogarithms. In particular a separated divergent endpoint value is
never silently treated as an ordinary integral. The same algorithm with
a variable upper endpoint gives an antiderivative; its value at zero
fixes the anchoring constant. Dominated convergence in the original
integral gives the asserted coincident-scale limits.
\end{proof}

\begin{remark}[What this settles]
For two positive integer orders, the theorem answers the finite-closure
part of the repository question on its entire integer product strip,
including $a=-1$. It gives an upper bound on representation complexity,
not a lower bound on arithmetic depth. Nor does it turn derivatives in
a noninteger spectral order into ordinary finite-weight hyperlogarithms.
\end{remark}
